\documentclass[11pt]{article}
\usepackage[margin=1in]{geometry}
\usepackage{parskip}
\usepackage{graphicx}

\title{EECS 473 Homework 0}
\author{Kevin Jin}
\date{Sep 4, 2026}

\begin{document}
\maketitle

\section*{Proposal 1: Through-Wall Motion Radar}

\subsection*{High-Level Description}

A portable radar that shows you people moving on the other side of a wall. The finished device is a box you hold in one hand (may be unrealistic, backup is two hands). Ideally we'd get the box to be around the size of a thick phone, but that may be unrealistic based on the antenna arrangement we will need to adopt. A patch antenna array points to the "top" of the box, a screen on the front showing a forward sector view. You point it at a wall, press a button, and moving people behind that wall appear as blips that track as they walk around. 

Commercial through-wall radars exist. The Range-R unit used by US police and the Xaver line are two examples, but they cost between six and fifty thousand dollars. The pitch is that a student team can build a functional one for a few hundred dollars in parts (that's with our designs and RnD cost excluded). The users who care are first responders searching the ruins of a collapsed building or tactical teams clearing a structure, both teams would want to know whether a room is occupied before they walk into it. 

The device works as an FMCW (frequency-modulated continuous-wave) radar that the team build themselves, operating somewhere in the 3 to 7 GHz range. The transmitter sends a chirp, the receiver mixes the return against the outgoing signal, and the beat frequency of that mix is proportional to target range. Multiple receive antennas give us angle of arrival by comparing phase across the array. Walls and furniture are static, so we subtract a running background estimate, a technique called moving-target indication, and what remains is the moving people. We render those returns as blips on the display.

% if using, put in another section
% We chose this frequency band deliberately. The 60 GHz mmWave radar modules that are cheap and available off the shelf do not penetrate walls, because attenuation at that frequency is severe. Lower frequencies pass through drywall, plywood, and brick with a manageable loss, and 3 to 7 GHz is the range where enough bandwidth is still available for useful range resolution. MIT's WiTrack project demonstrated through-wall localization of a single moving person in roughly this band with about 1.8 GHz of bandwidth, so the underlying physics is already proven. We are re-implementing a known result on a budget, not inventing a new one.

\subsection*{Implementation Issues}

\textbf{Rough plan for the Design Expo.} Allowing for the roughly ten-day lag between ordering a PCB and receiving it, and for the fact that RF boards usually require more than one revision, the schedule is approximately:

\begin{enumerate}
\item Planning and part selection (weeks 1 to 3). Choose the PLL/VCO, mixer, LNA, ADC, and processor. Simulate the antenna array. Settle on chirp bandwidth and chirp rate.
\item Breadboard the RF chain on evaluation modules (weeks 3 to 6). One transmit channel, one receive channel. Obtain a clean range measurement in free space and stream raw IQ samples to a laptop.
\item PCB design and first order (weeks 5 to 8). RF front end, processor, and power on a controlled-impedance board, with footprints for the full receive array. Order by roughly October 26.
\item DSP pipeline on the laptop (weeks 6 to 10). FMCW range processing, moving-target indication, angle of arrival, and the forward-sector visualization. Validate it through a real interior wall.
\item Board bring-up and integration (weeks 10 to 12). Move the DSP onto the processor, or stream to a companion module if necessary. Integrate the screen and enclosure. Order a second board revision if needed.
\item Final testing, calibration, poster, and Expo (weeks 12 to 14). Characterize range and accuracy, and add stretch goals if time allows.
\end{enumerate}

\textbf{Rough budget.} Across a team of five the ceiling is \$1250. The estimate below uses single-unit Digi-Key pricing for the RF semiconductors and current prototype-quantity PCB pricing. For calibration, a hobbyist chip-level FMCW RF board (VCO, LNA, mixer, three gain blocks) has been built for about \$60, and the classic MIT ``coffee-can'' radar runs a few hundred dollars in parts; our cost is higher mainly because we run multiple receive channels and fabricate a final board on Rogers laminate.

\begin{center}
\small
\begin{tabular}{p{6.5cm} p{5.5cm} c}
\hline
Category & Examples & Est. \$ \\
\hline
RF semiconductors & VCO, PLL, LNAs, mixers, gain blocks & 250 \\
ADC and digital capture & dual 12-bit 40--65 MSPS parts & 100 \\
Processor / FPGA board & Zynq-7020 class & 150 \\
RF connectors, cable, passives & SMA launches, RG405, filters, baluns & 150 \\
PCB fabrication & FR-4 iteration plus final Rogers hybrid & 300 \\
Display, battery, enclosure & 5 in LCD, LiPo, 3D-printed housing & 90 \\
Shipping & several part and board orders & 75 \\
\hline
\textbf{Total} & & \textbf{1115} \\
\hline
\end{tabular}
\end{center}

That leaves roughly \$135 of headroom under the \$1250 ceiling. Substituting a department-owned FPGA board for the Zynq purchase drops the total to about \$965. The Rogers PCB is the cost driver and the item most likely to cause schedule problems, since every revision is another board order and another ten-day wait, so we keep early iteration on cheap FR-4 and commit to the Rogers laminate only once the layout is stable.

\textbf{Other implementation issues.} RF design is unforgiving, and the antenna and front end together are the critical path for the whole project, so we are budgeting for at least two board revisions on their account. We need bench access to a vector network analyzer and a spectrum analyzer for bring-up, which the department lab provides. On the regulatory side, we stay within the FCC Part 15 UWB emission mask: transmit power stays very low, and no license is required. Building a high-power radar in a licensed band is not an option we are considering. There are no moving parts and no mechanical hazards. The energy budget is not a concern, because the transmit power is small and the display and processor are the dominant loads.

\textbf{Why we think it is doable.} FMCW radar in this frequency range is well-documented territory. MIT's 6.S062 ``coffee-can radar'' and Henrik Forsten's open 6 GHz FMCW radar project are both public end-to-end designs, and integrated PLL/VCO parts such as the ADF4159, paired with the HMC series, are available at reasonable cost. The hard and novel work for us is the antenna array, the moving-target indication through a real wall, and getting the pipeline to run in real time. None of those is an open research question. They are all work, and this team can do a large amount of work.

\subsection*{Evaluation and Testing}

The minimum measure of success is to detect and range a single person walking behind one interior drywall wall, with range error under about 50 cm. The next target is to add angle of arrival, so that the device tracks that person's position within the forward sector and the blip moves on the screen in real time as they walk.

The test method is as follows. A person walks a marked path in the far room while a camera in that room, timestamp-synchronized to the radar, records ground truth. We compare the estimated track against the actual track and extract detection range, range accuracy, angular accuracy, update rate, and the number of people we can resolve at once.

We will be honest in the report about the real-world caveats. The device only detects moving people, because the static subtraction erases anyone holding perfectly still; a breathing-detection mode is a stretch goal that would recover some of them. Wall material changes the propagation speed, so range requires a per-wall calibration adjustment, and we will not pretend the device works without one. Multiple walls and exterior brick are out of scope. Full imaging, meaning an actual silhouette or skeleton, is explicitly not what this project is; that would take years of work and a trained neural network. We track blips.

\newpage

\section*{Proposal 2 (Unrealistic): Head-Mounted Aurora Viewer}

\subsection*{High-Level Description}

A visor that lets you see auroae the naked eye cannot. The human eye integrates light over roughly a tenth of a second, so a faint aurora appears as a dim grey smudge with no color. A camera can integrate for five seconds and pull out bright, structured green and purple curtains. This device performs that integration live and puts the result on a display in front of your eyes.

The finished device is a visor: a low-light camera at the top, two micro-OLED displays, an IMU on the frame, and a compute pack on the headband or belt. It captures frames continuously, uses the IMU together with optical flow to register each new frame against the running stack so that head motion cancels out, and integrates about five seconds of exposure. The stacked, denoised, color-boosted image is sent to the displays at a few frames per second. The users are aurora chasers and astro-tourists who have driven somewhere cold and want to actually see the display, along with photographers who want a live preview instead of checking the back of a DSLR.

\subsection*{Implementation Issues}

\textbf{Why this one is risky.} The hard part is real-time multi-frame registration and stacking within a wearable power budget, and that is a computer-vision and machine-learning problem rather than an embedded-systems one. It needs a Jetson-class module, careful low-light denoising, and solid motion estimation. If we lean on a commercial compute module for all of that, our PCB becomes a carrier board with an IMU and a display driver, which undersells the goal of designing a PCB around a processor. Aurora also only occurs on cold, clear nights with low light pollution and sufficient geomagnetic activity, so field testing means driving north of the city at 1 a.m. hoping for a Kp index of 4.

\textbf{Rough plan for the Design Expo.} Weeks 1 to 3: select the camera, compute module, IMU, and displays; build a laptop prototype of the stacking pipeline against recorded night-sky footage. Weeks 4 to 7: design the carrier PCB (power, IMU, camera interface, display drivers) and order it. Weeks 6 to 10: port the pipeline to the compute module and get it running at a few frames per second. Weeks 10 to 13: integrate the visor, run field tests on clear nights, tune denoising and exposure. Weeks 13 to 14: poster and Expo, with an indoor simulated-sky demo as a fallback since aurora cannot be scheduled.

\textbf{Rough budget.} Dominated by the compute module, and well under the \$1250 ceiling.

\begin{center}
\small
\begin{tabular}{p{6.5cm} p{5.5cm} c}
\hline
Category & Examples & Est. \$ \\
\hline
Compute module & Jetson Orin Nano class & 250 \\
Low-light camera and lens & back-illuminated sensor, fast lens & 120 \\
Displays & two micro-OLED panels & 150 \\
IMU, power, passives & IMU, regulators, connectors & 60 \\
Carrier PCB & two revisions & 60 \\
Optics, visor frame, battery, enclosure & headband pack and housing & 110 \\
\hline
\textbf{Total} & & \textbf{750} \\
\hline
\end{tabular}
\end{center}

\textbf{Other implementation issues.} The compute module runs hot and draws several watts, so the headband pack needs a real battery and some airflow. Optics and camera selection is finicky and iterative. There are no RF or mechanical hazards, and the only safety note is not walking into things while wearing a display.

\subsection*{Evaluation and Testing}

We would compare the stacked output's signal-to-noise ratio against a single frame, and against a tripod-mounted long-exposure DSLR shooting the same patch of sky. Success is pulling a Kp 3 aurora into clear, colored visibility when the naked eye sees only a faint grey smudge. As an interim indoor test, we project a dim, slowly varying synthetic aurora onto a screen at known low light levels and confirm the device recovers it while a bare camera frame does not.

\end{document}
